\chapter{Schlüsseltausch-Verfahren}\label{ch:schluesseltausch}
 \begin{figure}[H]
	\centering
	\small

	\newcommand{\closedlock}[2]{%
		\begin{scope}[shift={#1},scale=#2]
			% Bügel
			\draw[line width=.5pt]
			(-.22,.25) -- (-.22,.58)
			arc[start angle=180,end angle=0,radius=.22]
			-- (.22,.25);

			% Schlosskörper
			\draw[rounded corners=.035cm,line width=.5pt]
			(-.32,-.18) rectangle (.32,.30);

			% Schlüsselloch
			\fill (0,.09) circle (.045);
			\fill (-.045,-.12) -- (.045,-.12) -- (.025,.08) -- (-.025,.08) -- cycle;
		\end{scope}
	}

	\newcommand{\openlock}[2]{%
		\begin{scope}[shift={#1},scale=#2]
			% geöffneter Bügel
			\draw[line width=.5pt]
			(-.22,.25) -- (-.22,.58)
			arc[start angle=180,end angle=20,radius=.22];

			% Schlosskörper
			\draw[rounded corners=.035cm,line width=.5pt]
			(-.32,-.18) rectangle (.32,.30);

			% Schlüsselloch
			\fill (0,.09) circle (.045);
			\fill (-.045,-.12) -- (.045,-.12) -- (.025,.08) -- (-.025,.08) -- cycle;
		\end{scope}
	}

	\newcommand{\ring}[2]{%
		\draw[line width=.55pt] #1 circle [radius=#2];
	}

	\newcommand{\lockbox}[3]{%
		\begin{tikzpicture}[x=1cm,y=1cm]
			\ring{(.58,1.15)}{.27}
			\ring{(1.76,1.15)}{.27}

			\ifx#1c
				\closedlock{(.88,.58)}{.72}
			\else
				\openlock{(.88,.58)}{.72}
			\fi

			\ifx#2c
				\closedlock{(2.06,.58)}{.72}
			\else
				\openlock{(2.06,.58)}{.72}
			\fi

			#3
		\end{tikzpicture}
	}

	\newcommand{\xortable}[6]{%
		\small
		\begin{tabular}{rll}
			#1               & \texttt{#2} \\
			$\oplus\quad$ #3 & \texttt{#4} \\
			\hline
			#5               & \texttt{#6}
		\end{tabular}
	}

	\begin{tikzpicture}[scale=0.88, every node/.style={transform shape}]
		% Layout constants
		\def\colL{-2}        % left column x
		\def\colR{10}        % right column x

		% Spaltentitel
		\node[font=\bfseries] at (\colL, 0.4) {Absender};
		\node[font=\bfseries] at (\colR, 0.4) {Empfänger};

		% === SCHRITT 1: Absender schliesst & verschlüsselt ===
		\node[anchor=north, align=center] (lockA1) at (\colL, 0) {
			\lockbox{c}{o}{\draw[->,line width=.45pt] (.88,-.28) -- (.88,.18);}\\[.1em]
			{\small Absender schliesst}
		};
		\node[anchor=north, below=0.1cm of lockA1] (tabA1) {
			\xortable
			{\textcolor{RoyalBlue}{Klartext}}{\textcolor{RoyalBlue}{101011}}
			{\textcolor{red}{Absender-Schlüssel}}{\textcolor{red}{011011}}
			{Erster Krypto-Text}{110000}
		};

		% === SCHRITT 2: Empfänger schliesst & verschlüsselt ===
		\node[anchor=north, align=center] (lockB1) at (\colR, -1) {
			\lockbox{c}{c}{\draw[->,line width=.45pt] (2.06,-.28) -- (2.06,.18);}\\[.1em]
			{\small Empfänger schliesst}
		};
		\node[anchor=north, below=0.1cm of lockB1] (tabB1) {
			\xortable
			{Erster Krypto-Text}{110000}
			{\textcolor{ForestGreen}{Empfänger-Schlüssel}}{\textcolor{ForestGreen}{101010}}
			{Zweiter Krypto-Text}{011010}
		};

		% Pfeil 1: Erster Krypto-Text  (links → rechts)
		% Arrow from output row of tabA1 to input row of tabB1
		\coordinate (a1start) at ([xshift=0.8cm]$(tabA1.south east)!0.17!(tabA1.north east)$);
		\coordinate (a1end) at ([xshift=-0.8cm]$(tabB1.north west)!0.17!(tabB1.south west)$);
		\draw[->, line width=.65pt] (a1start) -- (a1end)
			node[pos=0.5, below, sloped, font=\small, align=center]
			{Erster Krypto-Text:\\ \texttt{110000}};

		% === SCHRITT 3: Absender öffnet & entschlüsselt ===
		\node[anchor=north, align=center] (lockA2) at (\colL, -5) {
			\lockbox{o}{c}{\draw[->,line width=.45pt] (.88,-.28) -- (.88,.18);}\\[.1em]
			{\small Absender öffnet}
		};
		\node[anchor=north, below=0.1cm of lockA2] (tabA2) {
			\xortable
			{Zweiter Krypto-Text}{011010}
			{\textcolor{red}{Absender-Schlüssel}}{\textcolor{red}{011011}}
			{Dritter Krypto-Text}{000001}
		};

		% Pfeil 2: Zweiter Krypto-Text  (rechts → links)
		% Arrow from output row of tabB1 to input row of tabA2
		\coordinate (a2start) at ([xshift=-0.8cm]$(tabB1.south west)!0.17!(tabB1.north west)$);
		\coordinate (a2end) at ([xshift=0.8cm]$(tabA2.north east)!0.17!(tabA2.south east)$);
		\draw[<-, line width=.65pt] (a2end) -- (a2start)
			node[pos=0.5, below, sloped, font=\small, align=center]
			{Zweiter Krypto-Text:\\ \texttt{011010}};

		% === SCHRITT 4: Empfänger öffnet & entschlüsselt ===
		\node[anchor=north, align=center] (lockB2) at (\colR, -7) {
			\lockbox{o}{o}{\draw[->,line width=.45pt] (2.06,-.28) -- (2.06,.18);}\\[.1em]
			{\small Empfänger öffnet}
		};
		\node[anchor=north, below=0.1cm of lockB2] (tabB2) {
			\xortable
			{Dritter Krypto-Text}{000001}
			{\textcolor{ForestGreen}{Empfänger-Schlüssel}}{\textcolor{ForestGreen}{101010}}
			{\textcolor{RoyalBlue}{Klartext}}{\textcolor{RoyalBlue}{101011}}
		};

		% Pfeil 3: Dritter Krypto-Text  (links → rechts)
		% Arrow from output row of tabA2 to input row of tabB2
		\coordinate (a3start) at ([xshift=0.8cm]$(tabA2.south east)!0.17!(tabA2.north east)$);
		\coordinate (a3end) at ([xshift=-0.8cm]$(tabB2.north west)!0.17!(tabB2.south west)$);
		\draw[->, line width=.65pt] (a3start) -- (a3end)
			node[pos=0.5, below, sloped, font=\small, align=center]
			{Dritter Krypto-Text:\\ \texttt{000001}};

		% --- Abschluss ---
		\node[font=\small, align=center, below=0.2cm of tabB2] {Empfänger erhält den Klartext};
	\end{tikzpicture}

	\caption{Schlüsseltausch mit öffentlichen Schlüsseln}
\end{figure}

\clearpage
% \begin{figure}[H]
% 	\centering

% 	\newcommand{\closedlock}[2]{%
% 		\begin{scope}[shift={#1},scale=#2]
% 			% Bügel
% 			\draw[line width=.5pt]
% 			(-.22,.25) -- (-.22,.58)
% 			arc[start angle=180,end angle=0,radius=.22]
% 			-- (.22,.25);

% 			% Schlosskörper
% 			\draw[rounded corners=.035cm,line width=.5pt]
% 			(-.32,-.18) rectangle (.32,.30);

% 			% Schlüsselloch
% 			\fill (0,.09) circle (.045);
% 			\fill (-.045,-.12) -- (.045,-.12) -- (.025,.08) -- (-.025,.08) -- cycle;
% 		\end{scope}
% 	}

% 	\newcommand{\openlock}[2]{%
% 		\begin{scope}[shift={#1},scale=#2]
% 			% geöffneter Bügel
% 			\draw[line width=.5pt]
% 			(-.22,.25) -- (-.22,.58)
% 			arc[start angle=180,end angle=20,radius=.22];

% 			% Schlosskörper
% 			\draw[rounded corners=.035cm,line width=.5pt]
% 			(-.32,-.18) rectangle (.32,.30);

% 			% Schlüsselloch
% 			\fill (0,.09) circle (.045);
% 			\fill (-.045,-.12) -- (.045,-.12) -- (.025,.08) -- (-.025,.08) -- cycle;
% 		\end{scope}
% 	}

% 	\newcommand{\ring}[2]{%
% 		\draw[line width=.55pt] #1 circle [radius=#2];
% 	}

% 	\newcommand{\lockbox}[3]{%
% 		\begin{tikzpicture}[x=1cm,y=1cm]
% 			\draw[gray,line width=.7pt] (0,0) rectangle (3.15,1.62);

% 			\ring{(.58,1.15)}{.27}
% 			\ring{(1.76,1.15)}{.27}

% 			\ifx#1c
% 				\closedlock{(.88,.58)}{.72}
% 			\else
% 				\openlock{(.88,.58)}{.72}
% 			\fi

% 			\ifx#2c
% 				\closedlock{(2.06,.58)}{.72}
% 			\else
% 				\openlock{(2.06,.58)}{.72}
% 			\fi

% 			#3
% 		\end{tikzpicture}
% 	}

% 	\newcommand{\xortable}[6]{%
% 		\begin{tabular}{r@{\hspace{1.2cm}}l}
% 			#1               & \texttt{#2} \\
% 			$\oplus\quad$ #3 & \texttt{#4} \\
% 			\midrule
% 			#5               & \texttt{#6}
% 		\end{tabular}
% 	}

% 	\begin{tabular}{c@{\hspace{2.2cm}}c}

% 		\begin{minipage}{3.8cm}
% 			\centering
% 			\lockbox{c}{o}{%
% 				\draw[->,line width=.45pt] (.88,-.28) -- (.88,.18);
% 			}

% 			\vspace{.25em}
% 			Absender schliesst
% 		\end{minipage}
% 		 &
% 		\xortable
% 		{\textcolor{RoyalBlue}{Klartext}}{\textcolor{RoyalBlue}{101011}}
% 		{\textcolor{red}{Absender-Schlüssel}}{\textcolor{red}{011011}}
% 		{Erster Krypto-Text}{110000}
% 		\\[2.2cm]

% 		\begin{minipage}{3.8cm}
% 			\centering
% 			\lockbox{c}{c}{%
% 				\draw[->,line width=.45pt] (2.06,-.28) -- (2.06,.18);
% 			}

% 			\vspace{.25em}
% 			Empfänger schliesst
% 		\end{minipage}
% 		 &
% 		\xortable
% 		{Erster Krypto-Text}{110000}
% 		{\textcolor{ForestGreen}{Empfänger-Schlüssel}}{\textcolor{ForestGreen}{101010}}
% 		{Zweiter Krypto-Text}{011010}
% 		\\[2.2cm]

% 		\begin{minipage}{3.8cm}
% 			\centering
% 			\lockbox{o}{c}{%
% 				\draw[->,line width=.45pt] (.88,-.28) -- (.88,.18);
% 			}

% 			\vspace{.25em}
% 			Absender öffnet
% 		\end{minipage}
% 		 &
% 		\xortable
% 		{Zweiter Krypto-Text}{011010}
% 		{\textcolor{red}{Absender-Schlüssel}}{\textcolor{red}{011011}}
% 		{Dritter Krypto-Text}{000001}
% 		\\[2.2cm]

% 		\begin{minipage}{3.8cm}
% 			\centering
% 			\lockbox{o}{o}{%
% 				\draw[->,line width=.45pt] (2.06,-.28) -- (2.06,.18);
% 			}

% 			\vspace{.25em}
% 			Empfänger öffnet
% 		\end{minipage}
% 		 &
% 		\xortable
% 		{Dritter Krypto-Text}{000001}
% 		{\textcolor{ForestGreen}{Empfänger-Schlüssel}}{\textcolor{ForestGreen}{101010}}
% 		{\textcolor{RoyalBlue}{Klartext}}{\textcolor{RoyalBlue}{101011}}
% 	\end{tabular}

% 	\caption{Schlüsseltausch}
% \end{figure}

\section{Drei-Wege-Schlüsseltausch}
\begin{myprocedure}[Kommunikationsprotokoll Three-Pass Protocol]\label{myprocedure:Komm}
	Das Kommunikationsprotokoll Three-Pass Protocol (auch Schlüsseltausch mit drei Durchgängen genannt) ermöglicht es Alice, einen geheimen Schlüssel $t$ an Bob zu senden, ohne dass sie zuvor einen gemeinsamen geheimen Schlüssel vereinbart haben. Dazu verwenden sowohl Alice als auch Bob jeweils einen eigenen zufälligen Schlüssel ($s_A$ bzw. $s_B$).

	\begin{description}
		\item[Ausgangssituation:] Alice besitzt einen zufälligen binären Schlüssel $s_A$ der Länge $n$. Bob besitzt ebenfalls einen zufälligen binären Schlüssel derselben Länge $n$.
		\item[Ziel:] Alice hat zuvor einen geheimen binären Schlüssel $t$ der Länge $n$ gewählt. Diesen Schlüssel (das ist hier der Klartext) möchte sie über einen unsicheren Kanal an Bob senden, ohne dass Unbefugte den Schlüssel erfahren.
	\end{description}
	\begin{enumerate}
		\item \textbf{A}lice verschlüsselt den zu verschickenden Schlüssel \textcolor{blue}{$t$} (Klartext) mit ihrem Schlüssel \textcolor{red}{$s_A$}:
		      \begin{align*}
			      k_A := \textcolor{blue}{t} \oplus \textcolor{red}{s_A}
		      \end{align*}
		      und sendet den entstandenen Kryptotext $k_A$ an Bob.
		\item \textbf{B}ob verschlüsselt die empfangene Nachricht $k_A$ nun auch mit seinem Schlüssel \textcolor{ForestGreen}{$s_B$}:
		      \begin{align*}
			      k_{AB} := k_A \oplus \textcolor{ForestGreen}{s_B}
		      \end{align*}
		      und sendet den Kryptotext $k_{AB}$ zurück an Alice.
		\item Alice entschlüsselt den Kryptotext $k_{AB}$ mit ihrem Schlüssel \textcolor{red}{$s_A$}:
		      \begin{align*}
			      k_B = k_{AB} \oplus \textcolor{red}{s_A}
		      \end{align*}
		      und sendet $k_B$ an Bob.
		\item Schliesslich entschlüsselt Bob den Kryptotext $k_B$ mit seinem Schlüssel \textcolor{ForestGreen}{$s_B$}:
		      \begin{align*}
			      \textcolor{blue}{t} = k_B \oplus \textcolor{ForestGreen}{s_B}
		      \end{align*}
		      und erhält dadurch den Klartext \textcolor{blue}{$t$}, welchen Alice ihm mitteilen möchte.
	\end{enumerate}
\end{myprocedure}
\noindent
Warum funktioniert das? Der Kern der Geschichte liegt darin, dass eine zweite Anwendung eines Schlüssels die erste Anwendung desselben Schlüssels löscht (rückgängig macht) und
zwar, auch wenn zwischen diesen zwei Anwendungen andere Schlüssel angewendet worden sind. Betrachten Sie die folgende Berechnung
\begin{align*}
	 & t\oplus s_A \oplus s_B\oplus s_A\oplus s_B =    \\
	 & t\oplus (s_A\oplus s_A)\oplus (s_B\oplus s_B) = \\
	 & t\oplus 0\oplus 0 = t.
\end{align*}

\begin{myexercise}
	Spielen Sie den Schlüsseltausch mit einer weiteren Person aus der Klasse durch.
	\begin{myanswer}
		\begin{itemize}
			\item Alice wählt den Schlüssel $s_A = 1011$ und den zu übermittelnden Schlüssel $t = 1100$.
			\item Bob wählt den Schlüssel $s_B = 0110$.
			\item Alice berechnet $k_A$ und sendet $k_A$ an Bob:
			      \[
				      \begin{array}{r}
					      1100         \\
					      \oplus\;1011 \\
					      \hline
					      0111
				      \end{array}
			      \]
			      Also $k_A = 0111$.
			\item Bob berechnet $k_{AB}$ und sendet $k_{AB}$ an Alice:
			      \[
				      \begin{array}{r}
					      0111         \\
					      \oplus\;0110 \\
					      \hline
					      0001
				      \end{array}
			      \]
			      Also $k_{AB} = 0001$.
			\item Alice berechnet $k_B$ und sendet $k_B$ an Bob:
			      \[
				      \begin{array}{r}
					      0001         \\
					      \oplus\;1011 \\
					      \hline
					      1010
				      \end{array}
			      \]
			      Also $k_B = 1010$.
			\item Bob berechnet den Klartext-Schlüssel $t$:
			      \[
				      \begin{array}{r}
					      1010         \\
					      \oplus\;0110 \\
					      \hline
					      1100
				      \end{array}
			      \]
			      Also $t = 1100$.
		\end{itemize}
	\end{myanswer}
\end{myexercise}


\subsection*{Sicherheit des Drei-Wege-Schlüsseltauschs}
Ist das Drei-Wege-Kommunikationsprotokoll sicher? Bietet es dieselbe Sicherheitsgarantie wie die Verwendung einer Truhe mit zwei Schlössern?

Wenn ein Kryptoanalyst nur einzelne Kryptotexte des Protokolls erhält und
das Verfahren nicht kennt, wirkt die Kommunikation als eine Folge von Zufallsbits und in diesem Sinne ist unsere Implementierung dieses Verfahrens sicher.  Wir müssen aber damit rechnen, dass der Angreifer das Kommunikationsprotokoll kennt (oder errät) und die zwei zufällig generierten Schlüssel das Einzige sind, was ihm unbekannt ist.

Wenn der Angreifer (Eve) alle drei Kryptotexte ($k_A, k_{AB}, k_B$) abfangen kann, kann er durch die folgenden Berechnungen den Klartext $t$ herausfinden:
\begin{align*}
	 & \textcolor{Red}{k_A} \oplus \textcolor{orange}{k_{AB}} \oplus \textcolor{yellow}{k_B} =                                                                                                                                 \\
	 & (\textcolor{Blue}{t} \oplus \textcolor{DarkGreen}{s_A}) \oplus (\textcolor{Red}{k_A} \oplus \textcolor{Purple}{s_B}) \oplus (\textcolor{orange}{k_{AB}} \oplus \textcolor{DarkGreen}{s_A}) =                            \\
	 & (\textcolor{Blue}{t} \oplus \textcolor{DarkGreen}{s_A}) \oplus (\textcolor{Red}{k_A} \oplus \textcolor{Purple}{s_B}) \oplus ((\textcolor{Red}{k_A} \oplus \textcolor{Purple}{s_B}) \oplus \textcolor{DarkGreen}{s_A}) = \\
	 & \textcolor{Blue}{t} \oplus (\textcolor{DarkGreen}{s_A} \oplus \textcolor{DarkGreen}{s_A}) \oplus (\textcolor{Purple}{s_B} \oplus \textcolor{Purple}{s_B}) \oplus (\textcolor{Red}{k_A} \oplus \textcolor{Red}{k_A}) =   \\
	 & \textcolor{Blue}{t} \oplus 0 \oplus 0 \oplus 0 =                                                                                                                                                                        \\
	 & \textcolor{Blue}{t}.
\end{align*}

\clearpage
\section{Diffie-Hellman-Merkle-Schlüsseltausch}

\begin{myremark}[Neue mathematische Operatoren]\label{rem:operatoren}
	In den nachfolgenden Protokollen und kryptographischen Verfahren werden zwei Notationen verwendet, die sich vom gewohnten Gleichheitszeichen ($=$) unterscheiden:
	\begin{itemize}
		\item \textbf{Definitions- und Zuweisungszeichen ($:=$):}\\
		      Das Symbol $:=$ bedeutet: \emph{wird definiert als} bzw. \emph{wird berechnet und zugewiesen}. Während ein einfaches Gleichheitszeichen ($=$) oft eine Gleichheit zwischen zwei bereits bekannten Ausdrücken feststellt (z.\,B.\ $2 + 2 = 4$), signalisiert $:=$ die erstmalige Festlegung oder Berechnung einer Variablen (z.\,B.\ $x := g^a \bmod p$, d.\,h.\ der Wert von $x$ wird neu durch diese Rechnung definiert).
		\item \textbf{Kongruenzzeichen ($\equiv$):}\\
		      In der Modulo-Arithmetik bedeutet die Schreibweise
		      \[
			      a \equiv b \pmod m
		      \]
		      dass die Zahlen $a$ und $b$ bei der Division durch $m$ \textbf{denselben Rest} ergeben (gelesen als: \emph{$a$ ist kongruent zu $b$ modulo $m$}).
		      
		      \textbf{Beispiel:} $9 \equiv 2 \pmod 7$, denn $9 = 1 \cdot 7 + \mathbf{2}$ und $2 = 0 \cdot 7 + \mathbf{2}$. In Rechenschritten darf eine Zahl daher jederzeit durch ihren kleineren Rest modulo $m$ ersetzt werden.
	\end{itemize}
\end{myremark}

\begin{myprocedure}[Protokoll \gls{DHM}]\label{myprocedure:Diffie}
	Das Kommunikationsprotokoll \gls{DHM} ermöglicht es Alice und Bob, gemeinsam über einen öffentlichen Kanal einen geheimen Schlüssel zu vereinbaren, ohne dass sie zuvor einen gemeinsamen geheimen Schlüssel ausgemacht haben. Im Gegensatz zum Drei-Wege-Schlüsseltausch basiert das \gls{DHM}-Protokoll auf der Schwierigkeit, das diskrete Logarithmusproblem zu lösen (siehe \cref{rem:dislog}), ein mathematisches Problem, zu dem es (bisher) keinen effizienten Lösungsalgorithmus gibt.
	\begin{description}
		\item[Ausgangssituation:] Alice und Bob haben sich zuvor öffentlich auf eine grosse Primzahl $p$ und eine positive natürliche Zahl $g$ geeinigt. Dabei ist $g$ kleiner als $p$.
		\item[Ziel:] Alice und Bob möchten gemeinsam mit einer öffentlichen Kommunikation einen Schlüssel $s_{AB}$ vereinbaren. Diesen Schlüssel darf keine Drittperson in Erfahrung bringen.
	\end{description}
	\begin{enumerate}
		\item Alice wählt zufällig eine positive ganze Zahl $a$ mit $a<p$ und hält diese geheim. Dann berechnet sie mit dieser geheimen Zahl:
		      \begin{align*}
			      x := g^a\mod p
		      \end{align*}
		      und sendet $x$ an Bob.
		\item Bob wählt zufällig eine positive ganze Zahl $b$ mit $b<p$ und hält diese geheim. Dann berechnet er die Zahl
		      \begin{align*}
			      y := g^b\mod p
		      \end{align*}
		      und sendet $y$ an Alice.
		\item Alice erhält $y$ von Bob und berechnet mit ihrer geheimen Zahl $a$ die Zahl
		      \begin{align*}
			      s_{AB} := y^a\mod p.
		      \end{align*}
		\item Bob berechnet mit dem erhaltenen $x$ und seiner geheimen Zahl $b$ die Zahl
		      \begin{align*}
			      s_{BA} := x^b\mod p.
		      \end{align*}
	\end{enumerate}
	Mithilfe der Rechengesetze der Modulo-Operation kann bewiesen werden, dass tatsächlich
	\begin{align*}
		s_{AB} = s_{BA}
	\end{align*}
	gilt und somit Alice und Bob dieselbe Zahl berechnet haben. Diese Zahl $s_{AB}$ ist der gemeinsame Schlüssel.
\end{myprocedure}

\begin{myexample}[Modulare Exponentiation durch wiederholtes Quadrieren]
	Warum ist die Berechnung von $g^a \pmod p$ selbst für grosse Zahlen so schnell durchführbar?
	
	Würde man $g^a$ zuerst als riesige Zahl ausmultiplizieren und erst ganz am Schluss $\pmod p$ rechnen, würde der Speicherplatz jedes Computers sofort gesprengt. Stattdessen nutzt man zwei Prinzipien:
	\begin{enumerate}
		\item \textbf{Wiederholtes Quadrieren (\emph{Square-and-Multiply}):} Man zerlegt den Exponenten in eine Summe von Zweierpotenzen.
		\item \textbf{Zwischenergebnisse sofort reduzieren:} Nach jeder Multiplikation wird sofort modulo $p$ gerechnet. Dadurch bleibt jedes Zwischenergebnis immer kleiner als $p$.
	\end{enumerate}

	\textbf{Beispiel:} Berechnen Sie $3^{13} \pmod 7$.
	
	Wir zerlegen den Exponenten $13$ in Zweierpotenzen: $13 = \mathbf{8} + \mathbf{4} + \mathbf{1}$. Nun berechnen wir schrittweise die benötigten Potenzen durch wiederholtes Quadrieren, wobei wir das Ergebnis der vorherigen Zeile direkt wiederverwenden:
	\begin{align*}
		3^1 & = 3                                                                       &  & \equiv \textcolor{red}{3} \pmod 7         \\
		3^2 & = 3^2 = 9                                                                 &  & \equiv \textcolor{blue}{2} \pmod 7        \\
		3^4 & = (\textcolor{blue}{3^2})^2 \equiv \textcolor{blue}{2}^2 = 4              &  & \equiv \textcolor{ForestGreen}{4} \pmod 7 \\
		3^8 & = (\textcolor{ForestGreen}{3^4})^2 \equiv \textcolor{ForestGreen}{4}^2 = 16 &  & \equiv \textcolor{purple}{2} \pmod 7
	\end{align*}
	Nun multiplizieren wir nur die passenden Potenzen zusammen, setzen die farbigen Zwischenergebnisse ein und reduzieren erneut:
	\[
		3^{13} = 3^8 \cdot 3^4 \cdot 3^1 \equiv \textcolor{purple}{2} \cdot \textcolor{ForestGreen}{4} \cdot \textcolor{red}{3} = 24 \equiv 3 \pmod 7.
	\]
	Selbst bei einem Exponenten mit hunderten Stellen sind auf diese Weise nur wenige tausend Rechenoperationen mit handlichen Zahlen nötig.
\end{myexample}

\begin{myexercise}
	Berechnen Sie $2^{11} \pmod{13}$ mithilfe des wiederholten Quadrierens (Zerlegung des Exponenten und Reduktion bei jedem Zwischenschritt).
	\begin{myanswer}
		\begin{enumerate}
			\item Exponent in Zweierpotenzen zerlegen: $11 = \mathbf{8} + \mathbf{2} + \mathbf{1}$.
			\item Potenzen durch wiederholtes Quadrieren berechnen:
			      \begin{align*}
				      2^1 & = 2                                                                       &  & \equiv \textcolor{red}{2} \pmod{13}         \\
				      2^2 & = 2^2 = 4                                                                 &  & \equiv \textcolor{blue}{4} \pmod{13}        \\
				      2^4 & = (\textcolor{blue}{2^2})^2 \equiv \textcolor{blue}{4}^2 = 16              &  & \equiv \textcolor{ForestGreen}{3} \pmod{13} \\
				      2^8 & = (\textcolor{ForestGreen}{2^4})^2 \equiv \textcolor{ForestGreen}{3}^2 = 9   &  & \equiv \textcolor{purple}{9} \pmod{13}
			      \end{align*}
			\item Zusammenmultiplizieren und reduzieren:
			      \[
				      2^{11} = 2^8 \cdot 2^2 \cdot 2^1 \equiv \textcolor{purple}{9} \cdot \textcolor{blue}{4} \cdot \textcolor{red}{2} = 72 \equiv 7 \pmod{13}
			      \]
			      (denn $72 = 5 \cdot 13 + 7$).
		\end{enumerate}
	\end{myanswer}
\end{myexercise}

\begin{myremark}[Diskretes Logarithmusproblem]\label{rem:dislog}
	Beim \gls{DHM}-Verfahren wird die Eigenschaft ausgenutzt, dass Potenzieren modulo einer Primzahl eine sogenannte \textbf{Einwegfunktion} darstellt:
	\begin{itemize}
		\item \textbf{Einfache Richtung:} Bei gegebenen Werten $g$, $a$ und $p$ lässt sich das Ergebnis
		      \[
			      x \equiv g^a \pmod p
		      \]
		      auch für sehr grosse Zahlen effizient und schnell berechnen (mittels modularer Exponentiation).
		\item \textbf{Schwere Richtung (Diskreter Logarithmus):} Kennt ein Angreifer lediglich die Basis $g$, die Primzahl $p$ und das Ergebnis $x$, so ist es extrem rechenaufwendig, den Exponenten $a$ zu bestimmen, für den
		      \[
			      g^a \equiv x \pmod p
		      \]
		      gilt. Man bezeichnet $a$ als den \emph{diskreten Logarithmus} von $x$ zur Basis $g$ modulo $p$.
	\end{itemize}
	Für ausreichend grosse Primzahlen $p$ (in der Praxis mindestens 2048 Bit) existiert bis heute kein effizienter Algorithmus, um dieses Problem in vertretbarer Zeit zu lösen.
\end{myremark}


\begin{myexercise}
	Führen Sie das \gls{DHM}-Protokoll mit einer weiteren Person aus der Klasse durch. Verwenden Sie
	\begin{align*}
		p := 13 \quad \text{und}\quad g := 2
	\end{align*}
	(Sie können auch eigene Werte wählen) als öffentlich bekannte Schlüssel.
	\begin{myanswer}
		\begin{itemize}
			\item Alice wählt $a = 9$ und berechnet $x = 2^9 \mod 13 = 5$. Sie sendet $x = 5$ an Bob.
			\item Bob wählt $b = 10$ und berechnet $y = 2^{10} \mod 13 = 10$. Er sendet $y = 10$ an Alice.
			\item Alice berechnet den gemeinsamen Schlüssel: $s_{AB} = y^a \mod p = 10^9 \mod 13 = 12$.
			\item Bob berechnet den gemeinsamen Schlüssel: $s_{BA} = x^b \mod p = 5^{10} \mod 13 = 12$.
		\end{itemize}
		Alice und Bob haben somit denselben Schlüssel $s_{AB} = s_{BA} = 12$ vereinbart.
	\end{myanswer}
\end{myexercise}

\begin{myexercise}
	Versuchen Sie das Kommunikationsprotokoll \gls{DHM} zu knacken, indem Sie
	für die gegebenen Werte $p, g, x$ und $y$ jeweils die geheimen Zahlen $a$ und $b$ berechnen und daraus
	dann den vereinbarten Schlüssel $s_{AB}$.
	\begin{aenum}
		\item $g = 3, p = 5, x = 4, y = 2$
		\item $g = 2, p = 13, x = 6, y = 11$
	\end{aenum}
	\begin{myanswer}
		\begin{aenum}
			\item $a = 2, b = 3$ und $s_{AB} = 2^{2}\mod 5 = s_{BA} = 4^3 \mod 5 = 4$
			\item $a = 5, b = 7$ und $s_{AB} = 11^5 \mod 13 = s_{BA} = 6^7 \mod 13 = 7$
		\end{aenum}
	\end{myanswer}
\end{myexercise}

\subsection*{Sicherheit des \gls{DHM}-Protokolls}
Das \gls{DHM}-Protokoll ist sicher, weil es (bisher) keinen effizienten Algorithmus gibt, um das diskrete Logarithmusproblem zu lösen. Ein Kryptoanalyst (Eve) kann zwar die öffentlichen Werte $p, g, x$ und $y$ in Erfahrung bringen, aber um daraus den gemeinsamen Schlüssel $s_{AB}$ zu berechnen, müsste sie entweder die geheime Zahl $a$ von Alice oder die geheime Zahl $b$ von Bob bestimmen. Dies entspricht dem Lösen des diskreten Logarithmusproblems, was (bisher) als schwierig gilt.

Allerdings ist das DHM-Protokoll nicht gegen einen Man-in-the-Middle-Angriff geschützt. Ein Angreifer (Eve) könnte sich zwischen Alice und Bob schalten und so tun, als ob sie Alice wäre, wenn sie mit Bob kommuniziert, und umgekehrt. Dadurch könnte Eve sowohl mit Alice als auch mit Bob jeweils einen eigenen gemeinsamen Schlüssel vereinbaren und so die Kommunikation zwischen den beiden belauschen. Diese Schwäche kann durch die Verwendung von digitalen Signaturen behoben werden, welche die Authentizität der Kommunikationspartner sicherstellen. Ein Beispiel eines solchen Signaturverfahrens ist das \gls{RSA}-Kryptosystem (siehe \cref{subsec:rsa}).
